\documentclass{beamer}
\usepackage{amssymb}
\usepackage{amsmath}

\begin{document}

  \begin{frame}
    \textit{Propiedades de los Números Reales}
    \newline\newline

    Toda la teoría a desarrollar se basa en las propiedades fundamentales
    de los números reales:
    \newline
    \begin{itemize}[<+->]
      \item Todo subconjunto no vacío y acotado superiormente tiene supremo.
      \item Estructura de campo ordenado.
    \end{itemize}
  \end{frame}


  \begin{frame}
    \textit{Definición}: Sea $X \neq \emptyset$. Una función 
    $d: X \times X \to \mathbb{R}^+$ es una \textbf{métrica} en 
    $X$ si satisface para todo $x, y, z \in X$:
    \newline

    \begin{itemize}[<+->]
      \item $d(x, y) = 0 \iff x = y$.
      \item $d(x, y) = d(y, x)$.
      \item $d(x, y) \leq d(x, z) + d(z, y)$.
    \end{itemize}

    \pause
    \vspace{0.5em}
    El par $(X, d)$ se denomina \textbf{espacio métrico}.
  \end{frame}

  \begin{frame}
    \textit{Definición}: Sea $E$ un espacio vectorial real. 
    Una \textbf{norma} sobre $E$ es una función $\|\cdot\|: E \to \mathbb{R}$ 
    que satisface para todo $v, w \in E$ y $\lambda \in \mathbb{R}$:
    \newline

    \begin{itemize}[<+->]
      \item $\|v\| \geq 0$, y $\|v\| = 0 \iff v = 0$.
      \item $\|\lambda v\| = |\lambda|\|v\|$.
      \item $\|v + w\| \le \|v\| + \|w\|$.
    \end{itemize}

    \pause
    \vspace{0.5em}
    \textit{Teorema}: Toda norma $\|\cdot\|$ induce una 
    métrica en $E$ definida por $d(v, w) = \|v - w\|$.
  \end{frame}

  \begin{frame}
    \textit{Definición}: Sea $(X, d)$ un espacio métrico, $x_0 \in X$ y $r > 0$.
    \newline

    \begin{itemize}[<+->]
      \item Se define la \textbf{bola abierta} centrada en $x_0$ de radio $r$ como:
            \[ B(x_0, r) = \{x \in X: d(x, x_0) < r\} \]
      \item Un conjunto $U \subseteq X$ es \textbf{abierto} si
            para todo $x \in U$ existe $r > 0$ tal que $B(x, r) \subseteq U$.
      \item Un conjunto $F \subseteq X$ es \textbf{cerrado} si su 
            complemento $X \setminus F$ es abierto.
    \end{itemize}
  \end{frame}

  \begin{frame}
    \textit{Teorema (Propiedades de Abiertos y Cerrados)}:
    \newline

    \begin{itemize}[<+->]
      \item $\emptyset$ y $X$ son conjuntos abiertos y cerrados.
      \item La unión arbitraria de conjuntos abiertos es abierta.
      \item La intersección finita de conjuntos abiertos es abierta.
      \item La intersección arbitraria de conjuntos cerrados es cerrada.
      \item La unión finita de conjuntos cerrados es cerrada.
    \end{itemize}
  \end{frame}

  \begin{frame}
    \textit{Definición}: Sea $(X, d)$ un espacio métrico y $A \subseteq X$.
    \newline
    \begin{itemize}[<+->]
      \item Un punto $a \in A$ es \textbf{interior} a $A$
            si existe $r > 0$ tal que $B(a, r) \subseteq A$.
      \item El \textbf{interior} de $A$, denotado por
            $\operatorname{Int}(A)$, es el conjunto de todos los 
            puntos interiores de $A$.
    \end{itemize}

    \pause

    \vspace{0.5em}
    \textit{Teorema}: Para todo $A \subseteq X$:
    \[ \operatorname{Int}(A) = \bigcup_{\substack{U \subseteq A \\ U \text{ abierto}}} U \]
  \end{frame}

  \begin{frame}
    \textit{Definición}: Sea $(X, d)$ un espacio métrico y $A \subseteq X$.
    \newline

    \begin{itemize}[<+->]
      \item Un punto $x \in X$ es un \textbf{punto de acumulación} 
            de $A$ si para todo $\epsilon > 0$:
            \[ B(x, \epsilon) \cap (A \setminus \{x\}) \neq \emptyset \]
      \item El conjunto de todos los puntos de acumulación de $A$
            se denota por $A'$ y se llama \textbf{conjunto derivado}.
    \end{itemize}

    \pause

    \vspace{0.5em}
    \textit{Teorema}: Un conjunto $A \subseteq X$ es cerrado si y solo si $A' \subseteq A$.
  \end{frame}

  \begin{frame}
    \textit{Definición}: La \textbf{clausura} (o cierre) de 
                $A \subseteq X$, denotada por $\overline{A}$, se define como:
    \[ \overline{A} = A \cup A' \]

    \pause

    \textit{Teorema}: Para todo $A \subseteq X$:
    \newline
    \begin{itemize}[<+->]
      \item $x \in \overline{A} \iff \forall \epsilon > 0,
            \, B(x, \epsilon) \cap A \neq \emptyset$.
      \item $\overline{A}$ es un conjunto cerrado.
      \item $A$ es cerrado $\iff A = \overline{A}$.
      \item $\overline{A} = 
             \displaystyle\bigcap_{\substack{F \supseteq A \\ F \text{ cerrado}}} F$.
    \end{itemize}
  \end{frame}

  \begin{frame}
    \textit{Teorema (Identidades de Dualidad)}:
    \newline\newline
    Dado $(X, d)$ un espacio métrico y $A \subseteq X$,
    el interior y la clausura se relacionan mediante el complemento:
    \newline

    \begin{itemize}[<+->]
      \item $\operatorname{Int}(A) = X \setminus \overline{X \setminus A}$
      \item $\overline{A} = X \setminus \operatorname{Int}(X \setminus A)$
    \end{itemize}
  \end{frame}

  \begin{frame}
    \textit{Definición}: Sea $(X, d)$ un espacio métrico y $A \subseteq X$.
    \newline

    \begin{itemize}[<+->]
      \item La \textbf{frontera} de $A$, denotada por $\partial A$, se define como:
            \[ \partial A = \overline{A} \cap \overline{X \setminus A} \]
      \item Un conjunto $A$ es \textbf{denso} en $X$ si $\overline{A} = X$.
    \end{itemize}

    \pause
    \vspace{0.5em}
    \textit{Teorema}:
    \newline

    \begin{itemize}[<+->]
      \item $\overline{A} = \operatorname{Int}(A) \cup \partial A$.
      \item $A$ es abierto $\iff \partial A \subseteq X \setminus A$.
      \item $A$ es cerrado $\iff \partial A \subseteq A$.
    \end{itemize}
  \end{frame}

\end{document}
